% EM structure and teaching revision 20260928
\subsection{转动稳定性、固有磁矩与进动}
平面小线圈磁矩为 $\boldsymbol\mu=IS\boldsymbol n$，在匀强外场中净力为零而力矩为 $\boldsymbol M=\boldsymbol\mu\times\boldsymbol B$。若电流大小恒定，转轴固定且垂直于磁场、磁矩始终在垂直转轴的平面内，转角 $\theta$ 也就是磁矩与磁场的夹角，势能为 $U=-\mu B\cos\theta$。
平行取向附近 $U\simeq-\mu B+\mu B\theta^2/2$，转动惯量为 $J$ 时
\[
 J\ddot\theta=-\mu B\theta,\qquad \omega_0=\sqrt{\frac{\mu B}{J}}.
\]
反平行取向对应势能极大值，微小偏转会继续增大。如果电流并非由外源维持，而是在线圈运动中感生出来，则 $\mu$ 也会变化，不能直接照搬恒磁矩势能模型。
对于 $\boldsymbol\mu=\gamma\boldsymbol L$ 且角动量主要沿自转轴的快速转动体系，$\dd\boldsymbol L/\dd t=\gamma\boldsymbol L\times\boldsymbol B$ 给出进动。力矩垂直于角动量时主要改变方向，而非其大小；其近似前提是自转远快于进动，且外力矩不足以在一个自转周期内显著改变姿态。

